\begin{table*}[t]\centering\small\caption{ECE and NLL (mean over 5 seeds; lower is better, bold marks the lowest mean only, not a significant difference) per shift condition. r$\alpha$: rotation by $\alpha^\circ$; n$\sigma$: Gaussian noise; MCD: MC dropout; Ens: 5-member ensemble; TS: MLP with ID-fitted temperature.}\label{tab:eceNll}
\begin{tabular}{lcccccccc}
\toprule
Shift & ECE MLP & ECE TS & ECE MCD & ECE Ens & NLL MLP & NLL TS & NLL MCD & NLL Ens \\
\midrule
r0 & 0.018 & 0.016 & \textbf{0.016} & 0.017 & 0.117 & 0.104 & \textbf{0.083} & 0.111 \\
r10 & 0.104 & 0.070 & \textbf{0.042} & 0.078 & 0.745 & 0.614 & \textbf{0.491} & 0.645 \\
r20 & 0.226 & 0.161 & \textbf{0.149} & 0.208 & 1.954 & 1.450 & \textbf{1.370} & 1.687 \\
r30 & 0.457 & 0.388 & \textbf{0.361} & 0.415 & 4.747 & 3.365 & \textbf{3.098} & 3.987 \\
r45 & 0.816 & 0.785 & \textbf{0.752} & 0.803 & 13.431 & 9.490 & \textbf{8.639} & 12.257 \\
r60 & 0.827 & 0.803 & \textbf{0.792} & 0.820 & 15.266 & 11.101 & \textbf{10.442} & 13.874 \\
n0.25 & 0.105 & 0.067 & \textbf{0.043} & 0.083 & 0.763 & 0.587 & \textbf{0.479} & 0.667 \\
n0.5 & 0.392 & 0.335 & \textbf{0.280} & 0.343 & 4.109 & 2.903 & \textbf{2.402} & 3.467 \\
n0.75 & 0.584 & 0.524 & \textbf{0.474} & 0.522 & 7.746 & 5.468 & \textbf{4.525} & 6.597 \\
n1.0 & 0.674 & 0.618 & \textbf{0.555} & 0.593 & 10.145 & 7.230 & \textbf{5.993} & 8.689 \\
r30+n0.5 & 0.674 & 0.612 & \textbf{0.579} & 0.616 & 9.953 & 7.048 & \textbf{6.156} & 8.673 \\
\bottomrule
\end{tabular}\end{table*}
\begin{table*}[t]\centering\small\caption{Accuracy (higher is better) and Brier score (lower is better), mean over 5 seeds. TS has the same accuracy as MLP by construction.}\label{tab:accbrier}
\begin{tabular}{lcccccccc}
\toprule
Shift & Acc MLP & Acc TS & Acc MCD & Acc Ens & BRIER MLP & BRIER TS & BRIER MCD & BRIER Ens \\
\midrule
r0 & 0.972 & 0.972 & \textbf{0.977} & 0.973 & 0.047 & 0.047 & \textbf{0.037} & 0.045 \\
r10 & 0.812 & 0.812 & \textbf{0.839} & 0.825 & 0.287 & 0.271 & \textbf{0.227} & 0.257 \\
r20 & 0.650 & 0.650 & \textbf{0.657} & 0.653 & 0.562 & 0.520 & \textbf{0.502} & 0.527 \\
r30 & 0.400 & 0.400 & \textbf{0.412} & \textbf{0.412} & 1.012 & 0.935 & \textbf{0.899} & 0.949 \\
r45 & 0.132 & 0.132 & \textbf{0.139} & 0.128 & 1.657 & 1.606 & \textbf{1.559} & 1.636 \\
r60 & \textbf{0.139} & \textbf{0.139} & 0.136 & 0.131 & 1.674 & 1.636 & \textbf{1.616} & 1.660 \\
n0.25 & 0.839 & 0.839 & \textbf{0.861} & 0.852 & 0.258 & 0.241 & \textbf{0.205} & 0.236 \\
n0.5 & 0.496 & 0.496 & \textbf{0.516} & 0.500 & 0.867 & 0.812 & \textbf{0.749} & 0.813 \\
n0.75 & 0.301 & 0.301 & \textbf{0.303} & 0.301 & 1.244 & 1.171 & \textbf{1.097} & 1.163 \\
n1.0 & 0.217 & 0.217 & 0.225 & \textbf{0.226} & 1.407 & 1.331 & \textbf{1.246} & 1.307 \\
r30+n0.5 & 0.200 & 0.200 & 0.196 & \textbf{0.203} & 1.425 & 1.345 & \textbf{1.303} & 1.354 \\
\bottomrule
\end{tabular}\end{table*}
\begin{figure}[t]\centering\includegraphics[width=\columnwidth]{shift_curves.pdf}\caption{ECE (top) and NLL (bottom) against shift intensity for rotation (left) and noise (right); mean $\pm$ SD over 5 seeds. The r0 point is shared by both families.}\label{fig:curves}\end{figure}
\begin{table*}[t]\centering\small\caption{Temperature scaling versus the uncalibrated MLP per shift. $\Delta$ECE\% is the change of mean ECE from the MLP to TS relative to the MLP; $p$ is a paired $t$-test over 5 seeds.}\label{tab:tsmlp}
\begin{tabular}{lccccccc}
\toprule
Shift & ECE MLP & ECE TS & $\Delta$ECE\% & $p$ & NLL MLP & NLL TS & $p$ \\
\midrule
r0 & 0.018 & 0.016 & -12 & 0.268 & 0.117 & 0.104 & 0.094 \\
r10 & 0.104 & 0.070 & -33 & 0.018 & 0.745 & 0.614 & 0.009 \\
r20 & 0.226 & 0.161 & -28 & 0.017 & 1.954 & 1.450 & 0.012 \\
r30 & 0.457 & 0.388 & -15 & 0.016 & 4.747 & 3.365 & 0.009 \\
r45 & 0.816 & 0.785 & -4 & 0.038 & 13.431 & 9.490 & 0.011 \\
r60 & 0.827 & 0.803 & -3 & 0.099 & 15.266 & 11.101 & 0.012 \\
n0.25 & 0.105 & 0.067 & -36 & 0.016 & 0.763 & 0.587 & 0.020 \\
n0.5 & 0.392 & 0.335 & -14 & 0.017 & 4.109 & 2.903 & 0.016 \\
n0.75 & 0.584 & 0.524 & -10 & 0.012 & 7.746 & 5.468 & 0.011 \\
n1.0 & 0.674 & 0.618 & -8 & 0.017 & 10.145 & 7.230 & 0.010 \\
r30+n0.5 & 0.674 & 0.612 & -9 & 0.020 & 9.953 & 7.048 & 0.010 \\
\bottomrule
\end{tabular}\end{table*}

We go hypothesis by hypothesis. Per-seed standard deviations are in the generated tables (\texttt{results/tables/main.md}); with five seeds, all verdicts are descriptive.

\textbf{Overall picture.} Every system is accurate and well calibrated in distribution (the largest mean ECE of the four is \rhval{main/mlp/rot0/ece/mean:3}, for the MLP, in Table~\ref{tab:eceNll}) and becomes severely over-confident under shift: for the MLP at r60, accuracy is \rhval{main/mlp/rot60/accuracy/mean:3} (Table~\ref{tab:accbrier}) while mean confidence remains close to 1 (Fig.~\ref{fig:conf}). ECE, NLL and Brier all increase monotonically with intensity for each family (Fig.~\ref{fig:curves}). Accuracy is nearly identical across systems (differences of at most a few points, e.g.\ MCD \rhval{main/mc-dropout/rot10/accuracy/mean:3} vs MLP \rhval{main/mlp/rot10/accuracy/mean:3} at r10), so the differences below are mainly calibration, not accuracy.

\textbf{H1 (TS helps in distribution): not supported by the registered test.} At r0, TS lowers mean NLL from \rhval{main/mlp/rot0/nll/mean:3} to \rhval{main/mlp-+-temperature-scaling/rot0/nll/mean:3} and ECE from \rhval{main/mlp/rot0/ece/mean:3} to \rhval{main/mlp-+-temperature-scaling/rot0/ece/mean:3}, but the paired tests give $p=\rhval{cmp/abl_oracle/mlp/rot0/nll/paired_p:3}$ and $p=\rhval{cmp/abl_oracle/mlp/rot0/ece/paired_p:3}$ (Table~\ref{tab:tsmlp}, Table~\ref{tab:hyp}). The direction agrees with the hypothesis; five seeds on a \rhval{const/n_test}-image test set cannot show that the in-distribution gain is real, which is plausible because the MLP is already calibrated here (the fitted $T$ averages \rhval{main/mlp-+-temperature-scaling/rot0/temperature/mean:3}, Table~\ref{tab:oracle}).

\textbf{H2 (ID-fitted TS does not stay calibrated): supported for ECE; the NLL clause is refuted.} TS ECE is \rhval{main/mlp-+-temperature-scaling/rot60/ece/mean:3} at r60 and \rhval{main/mlp-+-temperature-scaling/noise1.0/ece/mean:3} at n1.0, versus \rhval{main/mlp-+-temperature-scaling/rot0/ece/mean:3} at r0: both exceed twice the r0 value, the registered criterion, by a wide margin (Table~\ref{tab:hyp}). Even at mild shift ECE rises from \rhval{main/mlp-+-temperature-scaling/rot0/ece/mean:3} to \rhval{main/mlp-+-temperature-scaling/rot10/ece/mean:3} at r10 and \rhval{main/mlp-+-temperature-scaling/noise0.25/ece/mean:3} at n0.25 (Table~\ref{tab:tsmlp}). The relative ECE change from the MLP to TS shrinks with rotation, from \rhval{cmp/abl_oracle/mlp/rot10/ece/rel_delta:pct0}\% at r10 to \rhval{cmp/abl_oracle/mlp/rot60/ece/rel_delta:pct0}\% at r60 ($p=\rhval{cmp/abl_oracle/mlp/rot60/ece/paired_p:3}$), so at the strongest rotation TS is indistinguishable from no calibration; for noise it is \rhval{cmp/abl_oracle/mlp/noise0.25/ece/rel_delta:pct0}\% at n0.25 and \rhval{cmp/abl_oracle/mlp/noise1.0/ece/rel_delta:pct0}\% at n1.0. The second part of H2 as registered (NLL gain over the MLP shrinks with intensity) is wrong: TS lowers NLL at every shifted condition (the largest paired $p$ is \rhval{cmp/abl_oracle/mlp/noise0.25/nll/paired_p:3}, at n0.25, Table~\ref{tab:tsmlp}) and the absolute gain grows (\rhval{main/mlp/rot0/nll/mean:3} to \rhval{main/mlp-+-temperature-scaling/rot0/nll/mean:3} at r0, \rhval{main/mlp/rot60/nll/mean:2} to \rhval{main/mlp-+-temperature-scaling/rot60/nll/mean:2} at r60 as in Table~\ref{tab:eceNll}). TS therefore reduces the damage under shift, but does not restore calibration; the reduction is largely a better NLL, not a calibrated one.

\textbf{H3 (ensemble best on NLL and Brier at severe shift): refuted.} At r60 and n1.0 MC dropout has the lowest NLL (\rhval{main/mc-dropout/rot60/nll/mean:3} and \rhval{main/mc-dropout/noise1.0/nll/mean:3}) and Brier score (\rhval{main/mc-dropout/rot60/brier/mean:3} and \rhval{main/mc-dropout/noise1.0/brier/mean:3}) of the four systems, while the ensemble has NLL \rhval{main/deep-ensemble-5/rot60/nll/mean:3} and \rhval{main/deep-ensemble-5/noise1.0/nll/mean:3} (Tables~\ref{tab:eceNll}, \ref{tab:accbrier}, and rows H3 of Table~\ref{tab:hyp}). The ensemble improves on the single MLP in NLL and Brier at every shift, but its NLL is higher than that of TS at every condition (r60: \rhval{main/deep-ensemble-5/rot60/nll/mean:3} vs \rhval{main/mlp-+-temperature-scaling/rot60/nll/mean:3}).

\textbf{H4 (ensemble ECE below MC dropout at severe shift): refuted.} The sign is the opposite: MCD ECE is lower at r60 (\rhval{main/mc-dropout/rot60/ece/mean:3} vs \rhval{main/deep-ensemble-5/rot60/ece/mean:3}, $p=\rhval{cmp/main/deep-ensemble-5/rot60/ece/paired_p:3}$) and n1.0 (\rhval{main/mc-dropout/noise1.0/ece/mean:3} vs \rhval{main/deep-ensemble-5/noise1.0/ece/mean:3}, paired $p=\rhval{cmp/main/deep-ensemble-5/noise1.0/ece/paired_p:sci1}$; Table~\ref{tab:hyp}). MC dropout has the lowest mean ECE and NLL in all eleven conditions (paired differences from the ensemble, listed in Table~\ref{tab:mcdens}, are significant at $\alpha=\rhval{const/alpha:2}$, uncorrected, for NLL in all 11, ECE in 9 and Brier in 7; the non-significant ECE cases are r0 and r60), although at r0 its ECE ties with TS to three decimals (Table~\ref{tab:eceNll}). This contradicts the common finding that ensembles are the stronger uncertainty method~\cite{ovadia2019can,lakshminarayanan2016simple}; at least part of our result is a regularisation effect of the dropout-trained network (Section~\ref{sec:ablations}), and the ensemble members here have no regularisation and are only 5 in number.

\textbf{Where the method fails.} (i) At r45 and r60 all four systems are at chance accuracy and mean ECE lies between \rhval{main/mc-dropout/rot45/ece/mean:2} and \rhval{main/mlp/rot60/ece/mean:2}: no system is usable and scalar rescaling cannot help (Table~\ref{tab:eceNll}). (ii) The combined shift r30+n0.5 gives ECE of \rhval{main/mlp-+-temperature-scaling/rot30_noise0.5/ece/mean:3} for TS, higher than either r30 (\rhval{main/mlp-+-temperature-scaling/rot30/ece/mean:3}) or n0.5 (\rhval{main/mlp-+-temperature-scaling/noise0.5/ece/mean:3}) alone. (iii) ID-fitted TS is never the best system on ECE or NLL outside r0 and is never well calibrated outside r0.
\begin{table}[t]\centering\small\caption{Registered hypothesis tests. A and B are means over 5 seeds; A$-$B and the paired $t$-test $p$ over the 5 seeds are statistics of \texttt{rh compare}. H2 compares the two means with the registered factor of two and H3 compares means only; in H3, A is MC dropout, the best system other than the ensemble (B). A $p$ that would print as zero is given in scientific notation.}\label{tab:hyp}
\resizebox{\columnwidth}{!}{\begin{tabular}{lllcccc}
\toprule
Hyp. & Quantity (A $-$ B) & Shift & A & B & A$-$B & $p$ \\
\midrule
H1 & NLL: TS $-$ MLP & r0 & \rhval{main/mlp-+-temperature-scaling/rot0/nll/mean:3} & \rhval{main/mlp/rot0/nll/mean:3} & \rhval{cmp/abl_oracle/mlp/rot0/nll/delta:3} & \rhval{cmp/abl_oracle/mlp/rot0/nll/paired_p:3} \\
H1 & ECE: TS $-$ MLP & r0 & \rhval{main/mlp-+-temperature-scaling/rot0/ece/mean:3} & \rhval{main/mlp/rot0/ece/mean:3} & \rhval{cmp/abl_oracle/mlp/rot0/ece/delta:3} & \rhval{cmp/abl_oracle/mlp/rot0/ece/paired_p:3} \\
H2 & TS ECE: shifted (A), r0 (B) & r60 & \rhval{main/mlp-+-temperature-scaling/rot60/ece/mean:3} & \rhval{main/mlp-+-temperature-scaling/rot0/ece/mean:3} & -- & -- \\
H2 & TS ECE: shifted (A), r0 (B) & n1.0 & \rhval{main/mlp-+-temperature-scaling/noise1.0/ece/mean:3} & \rhval{main/mlp-+-temperature-scaling/rot0/ece/mean:3} & -- & -- \\
H3 & NLL: MCD $-$ Ens & r60 & \rhval{main/mc-dropout/rot60/nll/mean:3} & \rhval{main/deep-ensemble-5/rot60/nll/mean:3} & \rhval{cmp/main/deep-ensemble-5/rot60/nll/delta:3} & -- \\
H3 & Brier: MCD $-$ Ens & r60 & \rhval{main/mc-dropout/rot60/brier/mean:3} & \rhval{main/deep-ensemble-5/rot60/brier/mean:3} & \rhval{cmp/main/deep-ensemble-5/rot60/brier/delta:3} & -- \\
H3 & NLL: MCD $-$ Ens & n1.0 & \rhval{main/mc-dropout/noise1.0/nll/mean:3} & \rhval{main/deep-ensemble-5/noise1.0/nll/mean:3} & \rhval{cmp/main/deep-ensemble-5/noise1.0/nll/delta:3} & -- \\
H3 & Brier: MCD $-$ Ens & n1.0 & \rhval{main/mc-dropout/noise1.0/brier/mean:3} & \rhval{main/deep-ensemble-5/noise1.0/brier/mean:3} & \rhval{cmp/main/deep-ensemble-5/noise1.0/brier/delta:3} & -- \\
H4 & ECE: MCD $-$ Ens & r60 & \rhval{main/mc-dropout/rot60/ece/mean:3} & \rhval{main/deep-ensemble-5/rot60/ece/mean:3} & \rhval{cmp/main/deep-ensemble-5/rot60/ece/delta:3} & \rhval{cmp/main/deep-ensemble-5/rot60/ece/paired_p:3} \\
H4 & ECE: MCD $-$ Ens & n1.0 & \rhval{main/mc-dropout/noise1.0/ece/mean:3} & \rhval{main/deep-ensemble-5/noise1.0/ece/mean:3} & \rhval{cmp/main/deep-ensemble-5/noise1.0/ece/delta:3} & \rhval{cmp/main/deep-ensemble-5/noise1.0/ece/paired_p:sci1} \\
H5 & ECE: TS $-$ Oracle & r45 & \rhval{main/mlp-+-temperature-scaling/rot45/ece/mean:3} & \rhval{abl_oracle/ts-oracle-shifted-val/rot45/ece/mean:3} & \rhval{cmp/abl_oracle/ts-oracle-shifted-val/rot45/ece/delta:3} & \rhval{cmp/abl_oracle/ts-oracle-shifted-val/rot45/ece/paired_p:sci1} \\
H5 & ECE: TS $-$ Oracle & r60 & \rhval{main/mlp-+-temperature-scaling/rot60/ece/mean:3} & \rhval{abl_oracle/ts-oracle-shifted-val/rot60/ece/mean:3} & \rhval{cmp/abl_oracle/ts-oracle-shifted-val/rot60/ece/delta:3} & \rhval{cmp/abl_oracle/ts-oracle-shifted-val/rot60/ece/paired_p:sci1} \\
H5 & ECE: TS $-$ Oracle & n0.75 & \rhval{main/mlp-+-temperature-scaling/noise0.75/ece/mean:3} & \rhval{abl_oracle/ts-oracle-shifted-val/noise0.75/ece/mean:3} & \rhval{cmp/abl_oracle/ts-oracle-shifted-val/noise0.75/ece/delta:3} & \rhval{cmp/abl_oracle/ts-oracle-shifted-val/noise0.75/ece/paired_p:sci1} \\
H5 & ECE: TS $-$ Oracle & n1.0 & \rhval{main/mlp-+-temperature-scaling/noise1.0/ece/mean:3} & \rhval{abl_oracle/ts-oracle-shifted-val/noise1.0/ece/mean:3} & \rhval{cmp/abl_oracle/ts-oracle-shifted-val/noise1.0/ece/delta:3} & \rhval{cmp/abl_oracle/ts-oracle-shifted-val/noise1.0/ece/paired_p:sci1} \\
\bottomrule
\end{tabular}
}\end{table}
\begin{table}[t]\centering\small\caption{MC dropout versus the 5-member ensemble per shift: difference of the means over 5 seeds ($\Delta$, MC dropout minus ensemble; negative favours MC dropout) and paired $t$-test $p$, uncorrected (scientific notation where four decimals would print as zero).}\label{tab:mcdens}
\resizebox{\columnwidth}{!}{\begin{tabular}{lcccccc}
\toprule
Shift & $\Delta$ECE & $p$ & $\Delta$NLL & $p$ & $\Delta$Brier & $p$ \\
\midrule
r0 & \rhval{cmp/main/deep-ensemble-5/rot0/ece/delta:3} & \rhval{cmp/main/deep-ensemble-5/rot0/ece/paired_p:4} & \rhval{cmp/main/deep-ensemble-5/rot0/nll/delta:3} & \rhval{cmp/main/deep-ensemble-5/rot0/nll/paired_p:4} & \rhval{cmp/main/deep-ensemble-5/rot0/brier/delta:3} & \rhval{cmp/main/deep-ensemble-5/rot0/brier/paired_p:4} \\
r10 & \rhval{cmp/main/deep-ensemble-5/rot10/ece/delta:3} & \rhval{cmp/main/deep-ensemble-5/rot10/ece/paired_p:4} & \rhval{cmp/main/deep-ensemble-5/rot10/nll/delta:3} & \rhval{cmp/main/deep-ensemble-5/rot10/nll/paired_p:4} & \rhval{cmp/main/deep-ensemble-5/rot10/brier/delta:3} & \rhval{cmp/main/deep-ensemble-5/rot10/brier/paired_p:4} \\
r20 & \rhval{cmp/main/deep-ensemble-5/rot20/ece/delta:3} & \rhval{cmp/main/deep-ensemble-5/rot20/ece/paired_p:4} & \rhval{cmp/main/deep-ensemble-5/rot20/nll/delta:3} & \rhval{cmp/main/deep-ensemble-5/rot20/nll/paired_p:4} & \rhval{cmp/main/deep-ensemble-5/rot20/brier/delta:3} & \rhval{cmp/main/deep-ensemble-5/rot20/brier/paired_p:4} \\
r30 & \rhval{cmp/main/deep-ensemble-5/rot30/ece/delta:3} & \rhval{cmp/main/deep-ensemble-5/rot30/ece/paired_p:4} & \rhval{cmp/main/deep-ensemble-5/rot30/nll/delta:3} & \rhval{cmp/main/deep-ensemble-5/rot30/nll/paired_p:4} & \rhval{cmp/main/deep-ensemble-5/rot30/brier/delta:3} & \rhval{cmp/main/deep-ensemble-5/rot30/brier/paired_p:4} \\
r45 & \rhval{cmp/main/deep-ensemble-5/rot45/ece/delta:3} & \rhval{cmp/main/deep-ensemble-5/rot45/ece/paired_p:4} & \rhval{cmp/main/deep-ensemble-5/rot45/nll/delta:3} & \rhval{cmp/main/deep-ensemble-5/rot45/nll/paired_p:4} & \rhval{cmp/main/deep-ensemble-5/rot45/brier/delta:3} & \rhval{cmp/main/deep-ensemble-5/rot45/brier/paired_p:4} \\
r60 & \rhval{cmp/main/deep-ensemble-5/rot60/ece/delta:3} & \rhval{cmp/main/deep-ensemble-5/rot60/ece/paired_p:4} & \rhval{cmp/main/deep-ensemble-5/rot60/nll/delta:3} & \rhval{cmp/main/deep-ensemble-5/rot60/nll/paired_p:4} & \rhval{cmp/main/deep-ensemble-5/rot60/brier/delta:3} & \rhval{cmp/main/deep-ensemble-5/rot60/brier/paired_p:4} \\
n0.25 & \rhval{cmp/main/deep-ensemble-5/noise0.25/ece/delta:3} & \rhval{cmp/main/deep-ensemble-5/noise0.25/ece/paired_p:4} & \rhval{cmp/main/deep-ensemble-5/noise0.25/nll/delta:3} & \rhval{cmp/main/deep-ensemble-5/noise0.25/nll/paired_p:4} & \rhval{cmp/main/deep-ensemble-5/noise0.25/brier/delta:3} & \rhval{cmp/main/deep-ensemble-5/noise0.25/brier/paired_p:4} \\
n0.5 & \rhval{cmp/main/deep-ensemble-5/noise0.5/ece/delta:3} & \rhval{cmp/main/deep-ensemble-5/noise0.5/ece/paired_p:4} & \rhval{cmp/main/deep-ensemble-5/noise0.5/nll/delta:3} & \rhval{cmp/main/deep-ensemble-5/noise0.5/nll/paired_p:4} & \rhval{cmp/main/deep-ensemble-5/noise0.5/brier/delta:3} & \rhval{cmp/main/deep-ensemble-5/noise0.5/brier/paired_p:4} \\
n0.75 & \rhval{cmp/main/deep-ensemble-5/noise0.75/ece/delta:3} & \rhval{cmp/main/deep-ensemble-5/noise0.75/ece/paired_p:4} & \rhval{cmp/main/deep-ensemble-5/noise0.75/nll/delta:3} & \rhval{cmp/main/deep-ensemble-5/noise0.75/nll/paired_p:4} & \rhval{cmp/main/deep-ensemble-5/noise0.75/brier/delta:3} & \rhval{cmp/main/deep-ensemble-5/noise0.75/brier/paired_p:4} \\
n1.0 & \rhval{cmp/main/deep-ensemble-5/noise1.0/ece/delta:3} & \rhval{cmp/main/deep-ensemble-5/noise1.0/ece/paired_p:4} & \rhval{cmp/main/deep-ensemble-5/noise1.0/nll/delta:3} & \rhval{cmp/main/deep-ensemble-5/noise1.0/nll/paired_p:4} & \rhval{cmp/main/deep-ensemble-5/noise1.0/brier/delta:3} & \rhval{cmp/main/deep-ensemble-5/noise1.0/brier/paired_p:4} \\
r30+n0.5 & \rhval{cmp/main/deep-ensemble-5/rot30_noise0.5/ece/delta:3} & \rhval{cmp/main/deep-ensemble-5/rot30_noise0.5/ece/paired_p:4} & \rhval{cmp/main/deep-ensemble-5/rot30_noise0.5/nll/delta:3} & \rhval{cmp/main/deep-ensemble-5/rot30_noise0.5/nll/paired_p:4} & \rhval{cmp/main/deep-ensemble-5/rot30_noise0.5/brier/delta:3} & \rhval{cmp/main/deep-ensemble-5/rot30_noise0.5/brier/paired_p:4} \\
\bottomrule
\end{tabular}
}\end{table}
\begin{figure}[t]\centering\includegraphics[width=\columnwidth]{conf_temp.pdf}\caption{Left: mean confidence against accuracy for each system and each of the 11 conditions (points on the dashed line are calibrated on average). Middle, right: ID-fitted and oracle temperature (mean $\pm$ SD over 5 seeds) for rotation and noise.}\label{fig:conf}\end{figure}
